\documentclass[11pt]{article}
\usepackage[margin=1in]{geometry}
\usepackage{amsmath,amssymb}
\usepackage{microtype}
\usepackage{xcolor}
\usepackage[hidelinks]{hyperref}

\definecolor{deepblue}{RGB}{28,63,110}

\title{A Measure-Zero Periodic Set Containing Every Rational Parameter}
\author{}
\date{}

\begin{document}
\maketitle

Here is a particularly sharp warning against cardinality and measure arguments.
It is a genuine analytic, time-reversible Hamiltonian brake system---the same
broad dynamical language as the Burrau problem.

For \(0<u<1\), consider the two-degree-of-freedom Hamiltonian
\[
  H_u(x,y,p_x,p_y)
  =\frac12\bigl(p_x^2+x^2\bigr)
   +\frac12\bigl(p_y^2+u^2y^2\bigr),
\]
and begin at the brake point
\[
  x(0)=y(0)=1,
  \qquad
  p_x(0)=p_y(0)=0.
\]
Hamilton's equations give
\[
  x(t)=\cos t,
  \qquad
  y(t)=\cos(ut),
\]
and
\[
  p_x(t)=-\sin t,
  \qquad
  p_y(t)=-u\sin(ut).
\]

\section*{The second-brake condition}

A second brake occurs precisely when both momenta vanish.  Thus there must be
positive integers \(m,n\) for which
\[
  t=n\pi,
  \qquad
  ut=m\pi.
\]
Eliminating \(t\) gives
\[
  u=\frac{m}{n}\in\mathbb{Q}.
\]
Conversely, if \(u=m/n\), then \(t=n\pi\) is a second-brake time.  Time
reversal about that brake point produces a periodic orbit.  Hence
\[
  \boxed{
    \text{a second brake exists exactly when }u\in\mathbb{Q}
  }
\]
and therefore
\[
  \boxed{
    \text{every rational parameter is periodic, while every irrational
    parameter is nonperiodic.}
  }
\]

\section*{Why this is surprising}

The periodic-parameter set \(\mathbb{Q}\cap(0,1)\)
\begin{itemize}
  \item is countable;
  \item has Lebesgue measure zero;
  \item contains no interval;
  \item is surrounded arbitrarily closely by nonperiodic parameters;
  \item nevertheless contains every rational parameter, including every
        rational Euclid parameter used to generate a Pythagorean triple.
\end{itemize}

For irrational \(u\), the trajectory winds densely on its invariant
two-torus.  For rational \(u\), the two frequencies are commensurable and the
trajectory closes.

\begin{center}
  \color{deepblue}
  \fbox{\parbox{0.86\textwidth}{\centering
  In this analytic Hamiltonian brake problem, almost every parameter is
  nonperiodic, yet every single rational parameter is periodic.}}
\end{center}

Thus neither ``periodic parameters have measure zero'' nor ``almost every
parameter is nonperiodic'' can settle the Pythagorean--Burrau conjecture.  If
the Burrau brake map behaved like this elementary oscillator family, the
conjecture would be maximally false even though its periodic parameters still
formed a measure-zero set.  An exact structural or arithmetic obstruction is
therefore indispensable.

\end{document}
